\section{Model and Training Specification}
\label{app:model}

This appendix states the deep-learning model in full. Section~\ref{sec:methods} describes what the
model is for; what follows is what it is.

\subsection{Input representation}

One training example is a window of a single runner's $T = 30$ most recent activities, ordered in
time, with the activity to be predicted last. Each activity is resampled to $R = 128$ points spaced
equally in distance rather than in time, so that a point index refers to the same fraction of the
run regardless of its duration. An example is therefore a tensor
\begin{equation}
  X \in \mathbb{R}^{T \times R \times C}, \qquad T = 30,\; R = 128,\; C = 11,
\end{equation}
with channels in fixed order: date, latitude, altitude, cadence, cumulative elevation, distance,
heart rate, power, activity type, outdoor WBGT, and pace. Pace occupies the final index by
construction, since the loss and the readout both address the target channel positionally.

Runners with fewer than 30 recorded activities are retained by padding the window at the front. Pad
slots are exactly zero in every channel, and a per-activity key mask excludes them from
cross-activity attention, so a short history is treated as short rather than as a sequence of
zero-valued runs. Establishing that contract required care: the date channel is centered by
subtracting the window maximum, which is applied after padding and would otherwise leave pad slots
with a nonzero date and defeat the mask. The centering is therefore recorded before the shift and
the pad slots reset after it.

Two channels are masked on the target activity only: pace, which is the prediction target, and
power, which is close to a deterministic function of pace and would otherwise leak it. Both remain
visible on all preceding activities. Weather channels are never masked, since race-day conditions
are known in advance from a forecast at the moment a prediction would be used.

Pace is smoothed along the distance axis with a trailing 10-point moving average. A centered filter
was used initially and is wrong for this purpose: zero-padding at the window edges biases the first
and last points toward zero, which reports the end of a run as faster than it was. Since the late
stages of a hot run are where the effect under study appears, the label was biased against the
hypothesis. The trailing filter costs a phase lag of roughly five points, uniform across all runs
and therefore incapable of confounding a weather comparison.

\subsection{Tokenization and embedding}

Every channel is discretized and embedded rather than consumed as a real number. Channel $c$ carries
a fixed edge vector $e^{(c)}$, and a value $v$ maps to the token index $\mathrm{digitize}(v,
e^{(c)})$, which indexes a learned table $E^{(c)} \in \mathbb{R}^{B_c \times d}$ with $d = 64$. Edges
are fixed at initialization: the index enters an integer lookup, so no gradient reaches the edges
and their placement is a design choice rather than a learned quantity.

Most channels use $B_c = 500$ bins. Placement is per-channel. Sensor channels use quantile edges;
pace uses uniform edges over $[0, 40]\,\mathrm{min\,km^{-1}}$, which matters because the loss reads
its bin grid from the pace row and uniform spacing makes that grid a linear pace axis. WBGT uses 100
quantile edges. The asymmetry is deliberate. For sensor channels, quantile placement concentrates
resolution in the dense middle and starves the sparse tails, where long runs and hard efforts carry
the fitness signal; substituting quantile edges for sensors cost 0.07 in $R^2$ in a direct test. For
WBGT the binding constraint is the number of hot runs available, which no bin placement can change,
so placing 100 bins where the observations actually are is the better use of them.

An absent sensor is a distinct token rather than a zero. Heart rate and power are frequently missing,
and mapping a missing reading to zero collides with a genuine reading of zero. Missing values survive
tokenization as \texttt{NaN} and map to index $B_c$, an extra row in the embedding table.

The channel embeddings at each point are summed and scaled, then combined with learned positional
embeddings over both axes and layer-normalized:
\begin{equation}
  z_{t,r} \;=\; \mathrm{LayerNorm}\!\left(\frac{1}{\sqrt{C}}\sum_{c=1}^{C} E^{(c)}\!\left[k^{(c)}_{t,r}\right]
                \;+\; P^{\mathrm{act}}_{t} \;+\; P^{\mathrm{pt}}_{r}\right).
\end{equation}
The $1/\sqrt{C}$ scale and the normalization are not cosmetic. A plain sum has standard deviation
growing as $\sqrt{C}$, which overwhelms the fixed-scale positional term and destabilizes the first
pre-normalization block. Introducing both raised held-out $R^2$ from 0.673 to 0.710 and lowered
error from 6.99\% to 6.60\% at a cost of 128 parameters.

\subsection{Factored attention}

Attention is factored over the two axes instead of applied to the flattened $T \times R = 3{,}840$
sequence, whose quadratic cost would be prohibitive. Each of the $L = 5$ layers applies attention
within an activity, across the $R$ points, and then attention across activities at matched point
index. The within-activity pass reads the shape of a single run; the cross-activity pass is where a
runner's recent training history, and therefore the fitness estimate and any conditioning context,
is assembled. The cross-activity block uses a low-rank projection, and the padding mask enters here
as a key mask so that padded slots cannot be attended.

Attention is bidirectional. No causal mask is applied along the point axis, since the whole of each
preceding activity is legitimately available and only the target activity's pace is withheld.

Hyperparameters: $L = 5$ layers, 8 heads, feed-forward width 256, embedding width 64, dropout 0.1.
Total parameters 5{,}950{,}251.

\subsection{Objective}

The model does not regress pace. It predicts a distribution over the 500 pace bins at each point and
is trained by cross-entropy against the bin containing the observed pace,
\begin{equation}
  \mathcal{L} \;=\; -\frac{1}{R}\sum_{r=1}^{R} \log \pi_{r}\!\left[\,b(y_{r})\,\right],
  \qquad b(y) = \mathrm{digitize}\!\left(y, e^{(\mathrm{pace})}\right),
\end{equation}
where $\pi_r$ is the predicted distribution at point $r$ of the target activity and $y_r$ the
observed pace there. The loss is computed on the target activity only.

A classification head over a continuous quantity is an unusual choice and is made for two reasons.
It imposes no parametric error distribution, so a skewed or multimodal pace at a given point is
representable; and it yields a predictive distribution at no extra cost, from which an interval
follows directly. Point predictions are read off as the expectation $\hat{y}_r = \sum_j \pi_{r,j}
m_j$ over bin midpoints $m_j$, and an activity-level prediction is the mean over the $R$ points.

Because the loss addresses the pace grid positionally, the bin edges must be stacked in numeric
channel order. Sorting them lexicographically, as a naive traversal of the parameter tree does,
places channel 9 after channel 10 and silently scores pace against another channel's grid. The
signature is diagnostic and worth recording: training loss falls normally while held-out error
remains pinned near 132\%, the value obtained by predicting the center of the bin range.

\subsection{Optimization}

Adam with global-norm gradient clipping at 1.0. The learning rate warms linearly from 0 to
$1\times10^{-4}$ over 1{,}000 steps, then follows a cosine decay to $1\times10^{-5}$ over 200{,}000
steps. Batch size 64, data-parallel across two GPUs. Training runs 200{,}000 steps, roughly 20 hours.

Checkpoints are written whenever held-out $R^2$ or error improves. These two criteria do not peak
together: $R^2$ peaks near 115{,}000 steps while error continues to improve to roughly 190{,}000.
Both checkpoints are therefore evaluated and reported, since selecting on $R^2$ alone overstates
error by about 0.2 percentage points. The artifacts in this paper come from the error-selected
checkpoint.

\section{Construction of the Heat Artifacts}
\label{app:curve}

Both heat artifacts are read off the trained model by counterfactual sweeps. For a runner's window,
the WBGT channel of the target activity is overwritten with a constant $w$ while every other channel
is held fixed, and the model is re-evaluated. The difference between two such evaluations is
attributable to WBGT alone, since nothing else in the input has changed. The sweep covers 5 to
30\,\textdegree C in 0.5\,\textdegree C steps; artifacts are truncated to $[10, 28]$, a 37-point
grid, because the evaluation panel contains no held-out target above 29.8\,\textdegree C and the
region beyond it cannot be checked.

Write $D_a(w)$ for runner $a$'s swept prediction. The runner's sensitivity is the percentage
slowdown at 28\,\textdegree C relative to 10\,\textdegree C,
\begin{equation}
  s_a \;=\; 100\left(\frac{D_a(28)}{D_a(10)} - 1\right),
\end{equation}
and the runner's normalized shape is $\phi_a(w) = \left(D_a(w)/D_a(10) - 1\right)\big/ (s_a/100)$,
which equals 0 at 10\,\textdegree C and 1 at 28\,\textdegree C by construction. Separating shape
from scale is what makes the median meaningful: runners differ mostly in the steepness of the
response, so averaging unnormalized curves would confound the two.

The population curve is $g(w) = \bar{s}\,\cdot\,\mathrm{med}_a\,\phi_a(w)$, the median shape rescaled
by the median sensitivity of the reference sample. The monotone projection $\max$-accumulate is
applied to the median shape at this point; as reported in Section~\ref{sec:methods} it moves 8 of 37
grid points by at most 0.13 percentage points.

A single window's sensitivity is dominated by noise, with test-retest reliability near 0.34, so
roughly three-quarters of the observed spread in $s_a$ across windows is measurement error rather
than a difference between runners. Averaging the 15 most recent windows raises reliability to
approximately 0.8. The per-runner score is therefore the percentile rank of a runner's window-averaged
sensitivity against the distribution of window-averaged sensitivities, not against the noise-inflated
single-window distribution. The prospective version uses exactly 15 preceding windows and excludes
the target activity.

\section{Compensability and the Convexity of the Response}
\label{app:physics}

This appendix supports the physical argument in Section~\ref{sec:results}. All coefficients are those
reported by \citet{taylor2006} and already used in Section~\ref{sec:literature}: convective exchange
$8.3\,(T_{sk}-T_a)\sqrt{v}$, radiant $5.2\,(T_{sk}-T_{mrt})$ and evaporative $124\,(P_{sk}-P_a)\sqrt{v}$
per square metre. Calculations assume a 70\,kg runner with $1.85\,\mathrm{m^2}$ of body surface at
three-and-a-half-hour marathon speed, $3.35\,\mathrm{m\,s^{-1}}$, with skin at $35\,^{\circ}$C.

\subsection{Wet-bulb temperature as a sufficient statistic}

For fully wetted skin the total capacity to shed heat is the sum of a dry and an evaporative term,
\begin{equation}
  Q_{\max} \;=\; h_c\,(T_{sk} - T_a) \;+\; h_e\,\bigl(P_s(T_{sk}) - P_a\bigr),
\end{equation}
with $P_s(\cdot)$ the saturation vapor pressure. The wet-bulb temperature is defined as the value a
wetted surface reaches when it exchanges no net heat with the air, so that $h_c T_a + h_e P_a = h_c
T_{wb} + h_e P_s(T_{wb})$. Substituting removes air temperature and humidity as separate quantities:
\begin{equation}
  Q_{\max} \;=\; h_c\,(T_{sk} - T_{wb}) \;+\; h_e\,\bigl(P_s(T_{sk}) - P_s(T_{wb})\bigr).
\end{equation}
This is an identity rather than an approximation. The ratio $h_e/h_c$ implied by the coefficients above
is $124/8.3 = 14.9\,^{\circ}\mathrm{C\,kPa^{-1}}$, whose reciprocal, $0.067\,\mathrm{kPa}\,^{\circ}
\mathrm{C}^{-1}$, is the standard sea-level psychrometric constant, so the two are mutually consistent.

Table~\ref{tab:isopleth} walks one wet-bulb isopleth. Across an 18-degree span of air temperature the
dry term falls from shedding $182\,\mathrm{W\,m^{-2}}$ to gaining $91$, the evaporative term compensates
exactly, and the total does not move.

\begin{table}[!ht]
\centering\small
\caption{Heat-loss capacity along the $22\,^{\circ}$C wet-bulb isopleth. Dry exchange and evaporation
trade off; their sum is invariant.}
\label{tab:isopleth}
\begin{tabular}{@{}rrrrr@{}}
\toprule
Air temp (\textdegree C) & Rel. humidity (\%) & Dry (W\,m$^{-2}$) & Evaporative & Total \\
\midrule
23 & 91.7 & $+182.3$ & 691.2 & 873.5 \\
29 & 54.3 & $+91.1$  & 782.4 & 873.5 \\
35 & 31.5 & $0.0$    & 873.5 & 873.5 \\
41 & 17.6 & $-91.1$  & 964.7 & 873.5 \\
\bottomrule
\end{tabular}
\end{table}

\subsection{Why the required fraction accelerates}

Thermal compensability is $w_{\mathrm{req}} = E_{\mathrm{req}}/E_{\max}$, where $E_{\mathrm{req}} =
H_{\mathrm{prod}} - (C + R)$ is the evaporation needed for balance and $E_{\max}$ the maximum the
environment will accept \citep{taylor2006}. Metabolic heat production for the runner specified above is
$424\,\mathrm{W\,m^{-2}}$. Table~\ref{tab:wreq} evaluates both parts. The numerator rises nearly
linearly; the denominator falls at an accelerating rate, because saturation vapor pressure gains 6.6\%
of its own value per degree at $12\,^{\circ}$C and 5.9\% at $26\,^{\circ}$C, on a base that has itself
nearly doubled. The ratio therefore accelerates: its increment per degree grows from 0.025 to 0.083
across the range.

\begin{table}[!ht]
\centering\small
\caption{Required fraction of maximal evaporation against wet-bulb temperature, at 60\% relative
humidity and no solar load.}
\label{tab:wreq}
\begin{tabular}{@{}rrrrr@{}}
\toprule
$T_{wb}$ (\textdegree C) & Air temp (\textdegree C) & $E_{\max}$ (W\,m$^{-2}$) & $w_{\mathrm{req}}$ & $\Delta w_{\mathrm{req}}$ per \textdegree C \\
\midrule
10.7 & 14.8 & 1047 & 0.012 & --- \\
15.0 & 19.8 & 962  & 0.118 & 0.0248 \\
19.3 & 24.8 & 851  & 0.253 & 0.0313 \\
23.7 & 29.8 & 704  & 0.453 & 0.0453 \\
28.1 & 34.9 & 516  & 0.817 & 0.0827 \\
\bottomrule
\end{tabular}
\end{table}

Comparing curvature over the same two bands used in Section~\ref{sec:results}, the learned curve steepens
by a factor of 2.12, from 0.31 to 0.66\% per degree. The compensability ratio steepens by 1.64 at 40\%
relative humidity with no sun, 1.85 at 70\% with $200\,\mathrm{W\,m^{-2}}$ of solar load, and 1.87 at
85\% with no sun. The prediction is of shape only: nothing in this calculation fixes the size of the
pace penalty, which depends on how much a runner concedes per unit of strain and is not derivable from
the heat balance.

\subsection{Why duration matters}

Balance need not hold instantaneously, since heat can be stored as a rise in core temperature. That
allowance is a rate, however, and it scales inversely with the duration over which it is spent. A
2.5\,\textdegree C rise in a 70\,kg runner represents about 607\,kJ, which is worth
$60{,}725\,\mathrm{W}$ over ten seconds but only $56\,\mathrm{W}$ over three hours, against a metabolic
production of $424\,\mathrm{W\,m^{-2}}$. A sprinter can therefore bank the entire thermal debt and never
pay it, which is consistent with \citet{guy2014} reporting faster sprint performances in the heat
alongside slower marathons, and it gives a reason beyond exposure time why the same conditions cost a
slower finisher more. \citet{taylor2006} list duration of competition among the determinants of
compensability for the same reason.

The calculations in this appendix are reproduced by \texttt{compute\_wettedness.py}.

\section{Baseline Specification}
\label{app:baseline}

The weather-blind baseline is gradient-boosted trees (XGBoost) on four feature groups: current-run
mean heart rate; trailing 30-day fitted pace; current-run distance; and elevation, entered as both
total gain and gain per kilometer. Settings: 400 trees, maximum depth 6, learning rate 0.04, subsample
0.8, column subsample 0.8, minimum child weight 20, histogram tree method. Three seeds (100, 200,
300) are averaged.

The target is $\log$ pace, with predictions returned to the pace scale by exponentiation. Fitting on
the log scale matches the multiplicative structure of the effect: heat is estimated and applied as a
percentage, so estimating it in the space in which it is applied avoids a scale mismatch. Note that
exponentiating a conditional mean of logs returns a conditional median rather than a mean, a
retransformation gap of order $\exp(\sigma^2/2)$. A smearing correction was tested and moved the
reported figures in the wrong direction, so no retransformation adjustment is applied; the level is
instead handled by the intercept in the calibration of Section~\ref{sec:methods}.

The supplementary branch replaces heart rate and trailing pace with a fixed, independently trained
marathon-fitness estimate together with distance and elevation. That baseline is available before a
run starts, whereas mean heart rate is not, so it tests whether the findings depend on conditioning
on a variable realized during the run being predicted.

\section{Architecture Search}
\label{app:ablation}

The architecture was not chosen a priori. Table~\ref{tab:ablations} records the modifications tested
against the accepted configuration of the moment, each as a separate 200{,}000-step run judged on the
same held-out panel of 1{,}349 windows from 158 runners. Changes were applied one at a time, since
bundling them makes attribution impossible.

\begin{table}[!ht]
\centering
\small
\caption{Architecture modifications tested. Each row was run against the accepted configuration at
the time it was tested, so entries are not all mutually comparable; the outcome column records the
decision. Unmarked rows are complete 200{,}000-step runs scored on the 1{,}349-window held-out panel.}
\label{tab:ablations}
\begin{tabular}{@{}llcc l@{}}
\toprule
Modification & Rationale & $R^2$ & Error (\%) & Outcome \\
\midrule
Trailing pace filter        & label bias at the finish   & 0.673 & 6.99 & kept \\
$1/\sqrt{C}$ scale $+$ norm & input scale grows with $C$ & 0.710 & 6.60 & kept \\
Padding mask $+$ missing token & pad and absence are not zero & 0.704 & 6.66 & kept \\
WBGT channel, 100 quantile bins & weather as input       & 0.710 & 6.61 & kept \\
Batch 32 $\rightarrow$ 64   & gradient variance          & 0.730 & 6.29 & kept \\[2pt]
Cross-block feed-forward restored & apparent missing capacity & 0.635 & 7.55 & rejected \\
Quantile bins for sensors   & bins wasted on outliers    & 0.635 & 7.08 & rejected \\
Window $T = 30 \rightarrow 60$ & longer history          & 0.693 & 6.83 & rejected \\
Live date channel$^{a}$     & recency is unencoded       & ---   & ---  & rejected \\
Random Fourier WBGT encoding & smooth encoding           & 0.709 & 6.52 & rejected \\
Embedding width 64 $\rightarrow$ 128$^{b}$ & capacity    & 0.802 & 6.74 & rejected \\
Acclimation channel         & carry exposure history     & 0.683 & 6.82 & rejected \\
\bottomrule
\end{tabular}
\end{table}

Several outcomes are informative beyond the selection they drove. Added capacity was consistently
harmful: restoring a feed-forward sublayer in the cross-activity block, and doubling embedding width,
both degraded every metric, which places the useful operating point well below the capacity the data
will support. Lengthening the history window also lost, and the diagnostic was that training loss was
worse as well, indicating a harder optimization rather than an overfitting gap; roughly half the
runners have fewer than 60 activities, so the additional slots are padding for them. Restoring the
dead date channel was the clearest failure, plausibly because slot position within a window is
already carried by the positional embedding and a second, noisier copy competes with it.

One negative result is methodological. The padding mask was implemented, smoke-tested against
synthetic all-zero slots, and shipped, but never fired in training: the date-centering step described
in Appendix~\ref{app:model} ran after padding and left every pad slot nonzero, so the detection test
that identified pad slots always returned false. The defect was silent, since a mask that never
activates produces a model that trains and evaluates normally. It was found only by testing the mask
on real pipeline output rather than on constructed input, and the accuracy gain previously credited
to the mask belongs to the missing-value token instead.

\section{Literature Conversions for Figure~\ref{fig:heateffects}}
\label{app:conversions}

Published heat effects are reported in incompatible units, and Figure~\ref{fig:heateffects} places them
on one axis. The conversions are as follows. Slopes given in seconds per degree are divided by a
representative finish time for the population studied, 130 minutes for elite fields and 240 minutes
for mass fields, giving a percentage per degree. Effects reported as a percentage over an interval,
such as a change per 5\,\textdegree C, are divided by the interval width. All curves are anchored at
10\,\textdegree C, so each shows slowdown relative to that reference rather than an absolute time.

Two limitations apply. Studies native to air temperature are drawn on the same axis as studies native
to WBGT, which are not interchangeable: the offset between them depends on humidity, radiation and
wind, so placement above roughly 18\,\textdegree C is approximate. And the published quadratics are
drawn only across the range over which they were fitted. A quadratic with a positive second-order
term diverges outside its support, and extrapolating one of the fitted clusters to 28\,\textdegree C
would imply a slowdown near 50\%, which is an artifact of the functional form rather than a claim
made by its authors.
\clearpage
\bibliographystyle{apalike}
\bibliography{references}
\clearpage

% Start of Appendix
\appendix
\renewcommand{\thetable}{A\arabic{table}}
\setcounter{table}{0} % Resets the table counter

% ===========================================================================
% v3 insert for Section 2, "How Heat Affects Endurance Running Performance".
% Placement: after Eq. (2) (the WBGT definition) and before "The exact
% functional form of the heat penalty remains contested" / Table 1.
%
% CHANGES FROM v2 (all verified by verify_heat_convexity_v3.py):
%  1. The flat-band paragraph is rebuilt. v2 attributed the ~16 C breakpoint to
%     the thermal ceiling descending to meet the aerobic one. That crossing is
%     at WBGT 25.7-27.1 C for this runner, ~11 C too warm to explain a 16 C
%     breakpoint, and v2 elsewhere rejects hard-ceiling behaviour. Both jobs are
%     now done by one convex strain-to-pace map, which is self-consistent.
%  2. Convexity is stated as a theorem, not a parameter sweep.
%  3. Curvature is framed as a LOWER bound (w_req alone gives 1.80; observed
%     2.18; a convex map lifts one to the other).
%  4. One coefficient convention: Taylor (2006) throughout, matching App. C.
%     This also makes the wet-bulb substitution exact to 0.01 C, so v2's
%     Lewis-relation footnote is no longer needed.
%  5. Robustness reported honestly: h_c and metabolic rate barely move the
%     prediction; T_sk dominates. The exposed-skin caveat is quantified.
%
% REQUIRES THREE EDITS ELSEWHERE IN THE MANUSCRIPT (see handoff notes):
%  * Section 2 currently uses T_sk = 35 C and says dry exchange reverses "near
%    36 C". Adopting 31 C moves the reversal to 31 C and changes the 83/14/4
%    evaporation shares. These must agree or a referee will catch it.
%  * Appendix C's isopleth table (\label{tab:isopleth}) is computed at
%    T_sk = 35 C, total 873.5. Table \ref{tab:wbtradeoff} below supersedes it at
%    T_sk = 31 C, total 555.1, so DELETE the appendix one rather than keeping
%    two isopleths at different skin temperatures. The labels differ, so
%    nothing breaks if you leave it in place -- it will just contradict.
%  * Appendix C reports the learned steepening as 2.12; this insert uses 2.18.
%    Settle which is current.
% ===========================================================================

\subsection{Why the penalty accelerates: a biophysical account}
\label{sec:convexity}

The forms collected in Table~\ref{tab:forms} are fitted rather than derived, and
the disagreement among them---linear, quadratic, piecewise---is a disagreement
about curvature. Standard heat-transfer relations predict curvature of a
particular sign, and the argument is short enough to give in full. Its
components are individually standard \citep{gagge1996, parsons2003, belding1955};
the assembly, and its use to predict the shape of a performance curve, we
believe to be new. \citet{jenkins2023} state the mechanism verbally---that the
skin-to-air vapour-pressure gradient becomes ``increasingly important to heat
balance and thus performance'' as conditions warm---without formalising it.

\paragraph{Wet-bulb temperature is a sufficient statistic.}
A runner whose skin is wetted rejects heat by two routes at once, dry exchange
with the air and evaporation of sweat:
\begin{equation}
  Q_{\max} = h_c\,(T_{sk} - T_a) \;+\; h_e\,\bigl(P_{s}(T_{sk}) - P_a\bigr),
  \label{eq:qmax-raw}
\end{equation}
with $h_c$ the convective heat-transfer coefficient, $h_e = \mathrm{LR}\cdot h_c$
its evaporative counterpart, $P_s(\cdot)$ saturation vapour pressure and $P_a$
ambient vapour pressure. Wet-bulb temperature is \emph{defined} as the
temperature a wetted surface reaches when it exchanges no net heat with the air,
so $h_c T_a + h_e P_a = h_c T_{wb} + h_e P_s(T_{wb})$. Substituting removes air
temperature and humidity as separate quantities:
\begin{equation}
  Q_{\max} = h_c\,\bigl(T_{sk} - T_{wb}\bigr)
           \;+\; h_e\,\bigl(P_{s}(T_{sk}) - P_{s}(T_{wb})\bigr).
  \label{eq:qmax-wb}
\end{equation}
The step is algebraic, and exact whenever the Lewis relation governing the
runner matches the one implicit in the definition of $T_{wb}$. The coefficients
of \citet{taylor2006} used throughout this paper give
$\mathrm{LR} = 124/8.3 = 14.9\,^{\circ}\mathrm{C\,kPa^{-1}}$, against a sea-level
psychrometric constant of $15.0$; the two definitions of $T_{wb}$ differ by
about $0.01\,^{\circ}$C at race conditions, so the identity holds to well within
any other uncertainty here.\footnote{Sources differ on this coefficient---ASHRAE
gives $16.5$ and ISO 7933 $16.7\,^{\circ}\mathrm{C\,kPa^{-1}}$ for human
sweating---and at $16.5$ the two wet-bulb definitions separate by about
$0.15\,^{\circ}$C. That does not disturb the argument, but \eqref{eq:qmax-wb} is
exact only under a single consistent choice.}
Table~\ref{tab:wbtradeoff} walks one such line. Every row sits at a wet-bulb
temperature of exactly $22\,^{\circ}$C and is therefore, by \eqref{eq:qmax-wb},
the same thermal environment---though the rows are assembled quite differently,
the first a cool morning at $92\%$ humidity and the last desert heat at $18\%$.
As air temperature rises the dry route degrades from shedding $121$ to
\emph{gaining} $152\,\mathrm{W\,m^{-2}}$, reversing sign at $31\,^{\circ}$C
where air and skin are equal, while the drier air widens the skin-to-air
vapour-pressure gradient and evaporative capacity rises by exactly as much. The
trade is exact rather than approximate: along a line of constant wet-bulb
temperature $h_e\,\mathrm{d}P_a = -h_c\,\mathrm{d}T_a$ by definition, so the two
terms move by equal and opposite amounts and their sum cannot shift. The
dominant weight on the wet-bulb term in \eqref{eq:wbgt} is therefore recoverable
from the heat balance rather than merely calibrated against outcomes.

\begin{table}[!ht]
\centering\small
\caption{One wet-bulb isopleth. Every row is at $T_{wb} = 22\,^{\circ}$C for the
reference runner of Appendix~\ref{app:physics}, a $70\,\mathrm{kg}$ athlete at
$5{:}00\,\mathrm{min\,km^{-1}}$ with skin at $31\,^{\circ}$C. Dry and
evaporative capacity trade off exactly, and their sum $Q_{\max}$ does not move.
The two right-hand columns show what that invariance does \emph{not} cover: the
strain the runner actually carries, and the fluid it costs, both rise along the
same line.}
\label{tab:wbtradeoff}
\begin{tabular}{@{}rrrrrrr@{}}
\toprule
 & & \multicolumn{3}{c}{Heat-loss capacity (W\,m$^{-2}$)}
   & \multicolumn{2}{c}{Cost to the runner} \\
\cmidrule(lr){3-5}\cmidrule(lr){6-7}
Air temp. & Rel.\ hum. & Dry & Evaporative & Total
   & Required & Sweat \\
(\textdegree C) & (\%) & $C$ & $E_{\max}$ & $Q_{\max}$
   & $w_{req}$ & (L\,h$^{-1}$) \\
\midrule
23 & 92 & $+121$ & 434 & 555 & 0.68 & 0.81 \\
29 & 54 & $+30$  & 525 & 555 & 0.74 & 1.06 \\
31 & 45 & $0$    & 555 & 555 & 0.75 & 1.14 \\
35 & 32 & $-61$  & 616 & 555 & 0.78 & 1.31 \\
41 & 18 & $-152$ & 707 & 555 & 0.80 & 1.56 \\
\bottomrule
\end{tabular}
\end{table}

\paragraph{Capacity falls concavely while requirement rises linearly.}
Two quantities move against one another as conditions warm, and they move at
different rates. The evaporation \emph{required} for balance,
$E_{req} = H_{prod} - (C + R)$, rises exactly linearly in air temperature,
because dry exchange is proportional to the skin-to-air gradient and simply
reverses sign once air passes skin temperature. The evaporative \emph{capacity}
falls, and falls at an accelerating rate, because saturation vapour pressure is
convex in temperature: by the Clausius--Clapeyron relation $\mathrm{d}P_s/
\mathrm{d}T$ rises from $0.09$ to $0.22\,\mathrm{kPa\,K^{-1}}$ between $12$ and
$28\,^{\circ}$C, so $P_s$ more than doubles across that span.\footnote{The
\emph{fractional} rate falls slightly over the same interval, from $6.6$ to
$5.8\%$ per kelvin. It is the absolute slope that governs the argument.}

The binding quantity is the ratio of the two, the heat stress index of
\citet{belding1955}, equivalently the skin wettedness required to stay in
balance:
\begin{equation}
  w_{req} \;=\; E_{req} / E_{\max}.
  \label{eq:wreq}
\end{equation}
That $w_{req}$ is convex follows from the two rates rather than from any choice
of parameter. Writing $w_{req} = N/D$, the numerator $N$ is affine and
increasing, and the denominator $D$ is positive, decreasing and concave; then
\begin{equation}
  w_{req}'' \;=\; \frac{-\,N D'' D \;-\; 2 N' D' D \;+\; 2 N D'^{2}}{D^{3}},
  \label{eq:convex}
\end{equation}
in which all three numerator terms are positive, since $D'' < 0$, $N' > 0$ and
$D' < 0$. Convexity is therefore a property of the heat balance and not an
artifact of the constants chosen. It is stated here on the air-temperature axis;
because wet-bulb globe temperature is very nearly affine in wet-bulb temperature
over the range plotted---its slope moves only from $1.047$ to $1.042$ across
$10$ to $38\,^{\circ}$C of air temperature---the property carries over to the
axis of Figure~\ref{fig:heat_cost}, which we confirm numerically at relative
humidities from 30 to 90\%.

Where $w_{req}$ would exceed the maximum sustainable wettedness $w_{\max}$, heat
stress is uncompensable at that speed: the ceiling caps sustainable metabolic
rate and the runner concedes pace until balance is restored. Nothing diverges,
and the consequence is a steep but finite decline.

Table~\ref{tab:wbtradeoff} also bounds how far the reduction of the previous
subsection extends. Wet-bulb temperature is a sufficient statistic for the
dissipation \emph{ceiling}, exactly, but not for the \emph{strain}. Because the
runner is a heat source of fixed output, $E_{req} = H_{prod} - C$ moves with the
dry term, so that
$w_{req} = (H_{prod} - C)/(Q_{\max} - C)$ is invariant only in the degenerate
case $H_{prod} = Q_{\max}$. Along the isopleth of Table~\ref{tab:wbtradeoff} it
instead climbs from $0.68$ to $0.80$, approaching ISO's unacclimatised limit of
$0.85$, while the sweat rate needed to hold pace nearly doubles from $0.81$ to
$1.56\,\mathrm{L\,h^{-1}}$. At matched wet-bulb, hot and dry is therefore harder
on a runner than cool and humid, because the dry route turning from ally to
adversary costs more than the widened vapour-pressure gradient returns. This
bears on the index rather than on the argument---the convexity result is
unaffected, and is in any case established below on the air-temperature
axis---but it is one principled reason to expect a model conditioned on WBGT
alone to leave residual variance.

\paragraph{One convex response explains both the flat band and the curvature.}
A proportional response to $w_{req}$ is not quite enough to describe the
measured curve. Rescaled to match the learned penalty at its endpoints,
$w_{req}$ predicts a $0.9\%$ slowdown at $16\,^{\circ}$C WBGT where the learned
curve shows $0.3\%$: strain rises perceptibly in the cool range while measured
pace does not. Nor can a hard ceiling supply the missing flatness. For the
reference runner the ceiling is reached only at WBGT $25.7$ to
$27.1\,^{\circ}$C, depending on whether $w_{\max}$ is taken as ISO's
unacclimatised $0.85$ or an acclimatised $1.0$, which is some eleven degrees too
warm to account for a breakpoint near $16\,^{\circ}$C.

What reconciles them is a single convex map from strain to conceded pace. Fitting
$\text{penalty} = c\,w_{req}^{\,k}$ over the $16$ to $28\,^{\circ}$C range
returns $k \approx 1.9$, and that one exponent does both jobs at once: it
suppresses the response at low strain, producing the flat band without a
threshold, and it steepens the response at high strain. This is physiologically
unsurprising---sweat produced beyond the evaporative ceiling drips rather than
cools, so the fluid cost compounds on exactly the days evaporation is least
effective \citep{taylor2006}---but it is fitted here rather than derived, and it
is the one link in the chain that the heat balance does not supply.

\paragraph{The duration dependence follows from the same balance.}
Balance need not hold instantaneously, since heat may be banked as a rise in
core temperature. That allowance is a \emph{rate}, $S = m\,c\,\Delta T_{core}/t$,
scaling as the inverse of duration: a tolerable $2.5\,^{\circ}$C rise gives a
$70\,\mathrm{kg}$ runner about $610\,\mathrm{kJ}$, which is $57\,\mathrm{W}$
spread over a three-hour marathon and $38\,\mathrm{W}$ over four and a half,
against roughly $770\,\mathrm{W}$ of heat production---about $7\%$ of the
problem at three hours, and less as the race lengthens. Over a ten-second sprint
the same allowance is worth tens of kilowatts, so heat balance is effectively
irrelevant at sprint distances and close to absolute at marathon distances.
This is the sign reversal \citet{guy2014} report across events, and it is a
second reason, beyond time on course, why the same conditions cost a slower
finisher more. It also predicts that the constraint should bind only weakly in
middle-distance work, which is consistent with \citet{leslie2024}, who find
critical velocity falling $3.7\%$ and then a further $2.9\%$ across heat indices
of $20$, $38$ and $55\,^{\circ}$C---a decelerating rather than accelerating
response, in precisely the regime where the storage term dominates.

\paragraph{Does it predict the curvature we measured?}
Between the $20$--$24$ and $24$--$28\,^{\circ}$C WBGT bands the curve learned in
Section~\ref{sec:results} steepens by a factor of $2.18$. Evaluated over the same
bands, $w_{req}$ alone steepens by $1.80$, and across relative humidities of
$40$--$85\%$ and solar loads of $0$--$500\,\mathrm{W\,m^{-2}}$ the predicted
factor spans $1.5$ to $2.7$. The prediction brackets the observation, and the
direction of any residual gap is the expected one: with a convex strain-to-pace
map the heat balance supplies a lower bound on curvature, not an equality.

The prediction is robust to most of its inputs and fragile in one. Varying the
convective coefficient across the full published span at this speed, $10.4$ to
$33.7\,\mathrm{W\,m^{-2}\,K^{-1}}$, moves the predicted factor only to
$1.4$--$2.8$; varying metabolic rate from a $60\,\mathrm{kg}$ runner at
$5{:}00\,\mathrm{min\,km^{-1}}$ to a $70\,\mathrm{kg}$ runner at $4{:}00$ moves
it to $1.5$--$2.8$. Mean skin temperature dominates instead. At the
laboratory-running value of $35\,^{\circ}$C the prediction falls to $1.3$--$1.8$
and misses the observation; at $29\,^{\circ}$C it widens to $1.6$--$4.3$. We
adopt $31\,^{\circ}$C because \citet{aylwin2023} measured mean skin temperature
falling from $29.8$ to $28.9\,^{\circ}$C through the men's marathon at the 2019
World Championships, below air temperature and well below laboratory treadmill
values, with running-speed airflow over wet skin the reason.

That measurement carries a caveat which cuts against us and is worth stating.
Aylwin's thermography covers exposed skin, some $64\%$ of body surface area,
while the heat balance wants a whole-body mean that includes covered skin, which
runs warmer. Weighting the covered remainder at anything from $32$ to
$38\,^{\circ}$C puts the whole-body mean between $30.3$ and $32.5\,^{\circ}$C,
over which range the predicted factor spans $1.4$ to $3.1$ and continues to
bracket the observation; the prediction fails only above about
$33\,^{\circ}$C, which no plausible weighting of these measurements reaches. The
agreement is therefore correct in sign and order of magnitude but not sharp, and
the width of the band is dominated by a quantity the corpus does not observe.

\paragraph{What this does and does not establish.}
The argument constrains shape, not magnitude: nothing here derives the $5.0\%$
penalty of Section~\ref{sec:results}. A model in which runners hold pace until
they strike the biophysical ceiling predicts almost no penalty until the high
$20$s and a collapse thereafter, which resembles no measured curve. Runners
evidently respond to graded strain long before the hard limit binds, and
$w_{req}$ is that graded quantity rather than a threshold.

Two competing explanations for the same curvature deserve naming, since neither
is excluded by our data. \citet{cheuvront2010} attribute an exponential decline
in aerobic performance to the cardiovascular axis: high skin temperature
compresses $\dot{V}\mathrm{O}_2$max reserve, so a fixed pace becomes a rising
relative intensity. \citet{wang2024} attribute more than half the marathon heat
effect to reduced atmospheric oxygen partial density, since warm air expands and
rising vapour pressure displaces oxygen---a mechanism that also scales with
vapour pressure, and therefore mimics the humidity signature of the evaporative
account while fitting a piecewise-\emph{linear} form. Distinguishing the three
requires measurements this corpus does not contain.

A statistical objection applies to any pooled curve. Because the penalty is
strongly ability-dependent \citep{ely2007}, differential slopes combined with
differential field composition across warm and cool races can bend an aggregate
curve while no individual curve is convex. Our design is partly protected: the
curve of Section~\ref{sec:results} is the median of \emph{per-runner}
counterfactual sweeps, in which each runner's own response is traced with every
other channel held fixed, rather than a regression pooled across runners and
races. It is not protected against the possibility that the deep model has
learned one population-average response and applies it uniformly.

Finally, the critical wet-bulb temperature at running metabolic rates appears
not to have been measured. \citet{wolf2021} report critical WBGT falling from
$35.3\,^{\circ}$C at rest to $31.6\,^{\circ}$C at $30\%$
$\dot{V}\mathrm{O}_2$max, about $163\,\mathrm{W\,m^{-2}}$; a marathon runs above
$400\,\mathrm{W\,m^{-2}}$. Occupational limit curves extrapolated that far give
$22\,^{\circ}$C WBGT (ACGIH) or $18\,^{\circ}$C (NIOSH), far below conditions
marathons are routinely held in, which suggests those curves do not transfer to
self-paced athletic work. The gap is worth naming as a target for measurement
rather than filled by extrapolation here.
